% Abhakseite „Das kann ich“ – Zone nur im Gesamt (\mitzone)
%% TODO Ich-kann-Sätze: Platzhalter wie in den Aufgabendateien
\begin{abhakseite}
\ifdefined\mitzone
\abhakgruppe{Kennst du schon}
\abhak{1}{Lösung durch Einsetzen prüfen, Ergebnis mit (wA)/(fA)}
\abhak{2}{Gleichung nach einer Variablen umstellen (Formel umstellen: ax + by = c nach y)}
\abhak{3}{Gerade aus der Gleichung zeichnen (n und Steigungsdreieck), Punkt ablesen, Lage zweier Geraden (parallel bei gleicher Steigung)}
\abhak{4}{Lineare Gleichung mit Klammern, Dezimalzahlen und x auf beiden Seiten lösen, Schreibform mit Strich}
\abhak{5}{Klammer ausmultiplizieren (Zahl mal Klammer, Minus vor der Klammer), Terme mit zwei Variablen zusammenfassen}
\abhak{6}{Rechnen mit Dezimalzahlen (Geld) und negativen Zahlen, Division mit dem Taschenrechner}
\abhak{7}{Gleichung mit einer Variablen aus einem Sachverhalt aufstellen (Preis · Anzahl, „zusammen“)}
\abhak{8}{Sek II: Terme mit drei Variablen zusammenfassen, Gleichungen vervielfachen und addieren, nach einer Variablen umstellen}
\abhak{9}{Sek II: lineare Ungleichungen in einer Variablen lösen und Lösungen sieben (Vorzeichenbedingung, Ganzzahligkeit, größter Wert)}
\abhak{10}{Sek II: Tripel als Punkte und Vektoren lesen und die Lösungsmenge als Menge schreiben}
\abhak{11}{Klammer ausmultiplizieren (Zahl mal Klammer, Minus vor der Klammer), Terme mit zwei Variablen zusammenfassen – Fehler finden}
\abhak{12}{Klammer ausmultiplizieren (Zahl mal Klammer, Minus vor der Klammer), Terme mit zwei Variablen zusammenfassen – selbst rechnen}
\fi
\abhakgruppe{Einheit 1 · Gleichungen mit zwei Variablen und grafisches Lösen}
\abhak{13}{Grafisch}
\abhak{14}{Grafisch, Sek II}
\abhak{15}{Grafisch – Fehler finden, Begründen}
\abhakgruppe{Einheit 2 · Einsetzungsverfahren}
\abhak{16}{Einsetzen}
\abhak{17}{Bruch als Lösung (Vorrat)}
\abhak{18}{Einsetzen – Fehler finden, Begründen}
\abhakgruppe{Einheit 3 · Additionsverfahren}
\abhak{19}{Addition}
\abhak{20}{Addition, Sek II}
\abhak{21}{Dezimalzahlen und Brüche}
\abhak{22}{Addition – Fehler finden, Begründen}
\abhakgruppe{Einheit 4 · Sachaufgaben}
\abhak{23}{Sachaufgaben}
\abhak{24}{Sachaufgaben, Sek II}
\abhak{25}{Mischungen und Alter (Vorrat)}
\abhak{26}{Probe im Text (Rechnung mit den Zahlen der Aufgabe)}
\abhak{27}{Bestandteile einer Gleichung deuten}
\abhak{28}{Sachaufgaben – Fehler finden, Begründen}
\abhakgruppe{Einheit 5 · Drei Variablen und Lösungsvielfalt (Sek II)}
\abhak{29}{Drei Variablen und Lösungsvielfalt}
\abhak{30}{Drei Variablen und Lösungsvielfalt – Fehler finden, Begründen}
\end{abhakseite}
